\input{../article_base.tex}
\title{פתרון מטלה 05 --- תבניות ריבועיות ומספרים P־אדיים, 80507}

\DeclareMathOperator{\Deg}{Deg}

% chktex-file 44

\begin{document}
\maketitle
\maketitleprint[red]

\question{}
יהי $\FF$ שדה, $\alpha \in \RR$ כך ש־$0 < \alpha < 1$.
נסמן $\FF((X^{-1}))$ ההשלמה של $\FF(X)$ ביחס לנורמה $| \cdot |_{\alpha}$ ואת $\FF[[X^{-1}]]$ חוג השלמים ב־$\FF((X^{-1}))$.

\subquestion{}
נראה שלכל $s \in \ZZ$ ולכל סדרה ${(a_n)}_{n = s}^{\infty} \subseteq \FF$ הטור,
\[
	\sum_{n = s}^\infty a_n X^{-n}
\]
מתכנס ב־$\FF((X^{-1}))$.
\begin{proof}
	אנו רוצים להראות התכנסות בהשלמה ולכן מספיק להראות שהסדרה ${(b_n)}_{n = s}^{\infty}$ המוגדרת על־ידי,
	\[
		\sum_{i = s}^n a_i X^{-i}
	\]
	היא סדרת קושי.
	נבחין כי אם $n > m \in \NN$ איברים גדולים מ־$s$ אז,
	\begin{align*}
		b_n - b_m
		& = \sum_{i = n}^m a_i X^{-i} \\
		\implies & |b_n - b_m|_{\alpha} \\
		& = \max\{ |a_i X^{-i}|_{\alpha} \mid m \le i \le n \} \\
		& = \max\{ \alpha^{-\Deg(a_i X^{-i})} \mid m \le i \le n \} \\
		& = \max\{ \alpha^{-(\deg(a_i X^{0}) - \deg(X^i))} \mid m \le i \le n \} \\
		& = \max\{ \alpha^{-(0 - i)} \mid m \le i \le n \} \\
		& = \alpha^m
	\end{align*}
	ולמעשה קיבלנו ישירות שזוהי סדרת קושי מאי־שוויון האולטרה־מטרי.
\end{proof}

\subquestion{}
נראה שלכל $\phi \in \FF[[X^{-1}]]$ ולכל $k \in \NN$ קיימים $a_0, \ldots, a_{k - 1} \in \FF$ יחידים כך שמתקיים,
\[
	| \phi - (a_0 + \cdots + \frac{a_{k - 1}}{X^{k - 1}}) |_{\alpha}
	\le \alpha^k
\]
\begin{proof}
	הוכחנו במטלה הקודמת שמתקיים עבור $(\FF(X), | \cdot |_{\alpha})$,
	\[
		O_{| \cdot |_{\alpha}}
		= \{ \frac{f}{g} \mid f, g \in \FF[X], g \ne 0, \deg f \le \deg g \}
	\]
	ובשיעור הראינו שחוג השלמים של ההשלמה הוא הסגור של חוג השלמים של השדה המקורי, לכן,
	\[
		\FF[[X^{-1}]]
		= \cl \{ \frac{f}{g} \mid f, g \in \FF[X], g \ne 0, \deg f \le \deg g \}
	\]
	בהתאם אם $h \in \FF[[X^{-1}]]$ אז קיימת סדרה $h_n \in O_{| \cdot |_{\alpha}}$ המקיימת,
	\[
		| h - h_n |_{\alpha}
		\to 0
	\]
	בפרט נוכל לבחור  $h_n$ עבור $\varepsilon = \alpha^k$.
	נבחין כי מאי־שוויון האולטרה־מטרי ומהלך מאוד דומה להוכחת טענה זו עבור שלמים $p$־אדיים נקבל שקיים פולינום,
	\[
		a_0 + \cdots + \frac{a_{k - 1}}{X^{k - 1}}
		\in \{ h_n \}
	\]
	המקיים את אי־השוויון, וזאת שכן המקדמים הגדולים יותר יתאפסו.
	לבסוף אם כך נותר להראות שפולינום זה הוא יחיד, גם הוכחה זו זהה לזו שהובאה בשיעור,
	נניח ש־${\{ b_n \}}_{n = 0}^{k - 1} \subseteq \FF$ סדרה המקיימת את הטענה ושונה מ־$a_0, \ldots, a_{k - 1}$.
	אז,
	\[
		\alpha^{\max\{ k \mid a_k \ne b_k \}}
		= |a_0 + \cdots + \frac{a_{k - 1}}{X^{k - 1}} - b_0 + \cdots + \frac{b_{k - 1}}{X^{k - 1}}|_{\alpha}
		\le \alpha^k
	\]
	וזוהי כמובן סתירה.
\end{proof}

\subquestion{}
נראה שלכל $\phi \in \FF[[X^{-1}]]$ קיימת סדרה ${\{ s_n \}}_{n = 0}^\infty \subseteq \FF$ כך שמתקיים,
\[
	\phi
	= \sum_{n = 0}^\infty \frac{a_n}{X^n}
\]
\begin{proof}
	מהסעיף הקודם נוכל להגדיר את $a_n$ להיות המקדם המוביל בפולינום היחיד המקיים את הטענה ל־$k = n$.
	בהתאם נגדיר את,
	\[
		f_k
		= \sum_{n = 0}^k \frac{a_n}{X^n}
	\]
	ונקבל סדרה כך ש־$|f_n - f_{n - 1}|_{\alpha} = |\frac{a_n}{X^n}|_{\alpha} = \alpha^n$ ולכן היא סדרת קושי ומתכנסת.
	לבסוף נקבל $|\phi - f_k|_{\alpha} \to 0$ ולכן היא סדרה המקיימת את הטענה בדיוק.
\end{proof}

\subquestion{}
נוכיח שלכל $\phi \in {\FF((X^{-1}))}^*$ קיימים $s \in \ZZ$ וסדרה ${(a_n)}_{n = s}^{\infty} \subseteq \FF$ כך ש־$a_s \ne 0$ וגם,
\[
	\phi = \sum_{n = s}^\infty a_n X^{-n}
\]
\begin{proof}
	קיים $s \in \ZZ$ יחיד המקיים,
	\[
		| \phi |_{\alpha}
		= \alpha^s
	\]
	ונגדיר $f = \phi X^{-s}$, נקבל ש־$|f|_{\alpha} = 1$, אם כך $f \in \FF[[X^{-1}]]$ ולכן קיימת לה סדרה יחידה ${(b_n)}_{n = 0}^{\infty} \subseteq \FF$ כך שמתקיים,
	\[
		f
		= \sum_{n = 0}^\infty \frac{b_n}{X^n}
	\]
	ולכן מהתכנסות בהחלט גם,
	\[
		\phi
		= \sum_{n = 0}^\infty \frac{b_n}{X^n} \cdot X^s
		= \sum_{n = s}^\infty \frac{b_{n + s}}{X^n}
	\]
	ונוכל להגדיר $a_n = b_{n + s}$.
	נשאר להראות ש־$a_s \ne 0$, אך אם זה המצב נוכל להגדיל את $s$ עד שנקבל $a_s \ne 0$, לא יתכן שתהליך זה לא יסתיים אחרת $\phi = 0$ בסתירה לשייכותו להפיכים.
\end{proof}

\question{}
יהי $\frac{m}{n}$ רציונלי, נתון כי,
\[
	\frac{m}{n}
	= {(\ldots 212121)}_3
\]
נמצא את $m$ ואת $n$.
\begin{solution}
	מהמחזוריות נקבל $r_l = r_{l + 2}$ לכל $l$.
	נסמן גם $r_0 = m$.
	אנו יודעים כי מתקיים,
	\[
		r_0 \equiv n a_0 \mod 3
		\iff m \equiv n \mod 3
	.\]
	נחשב,
	\[
		r_1
		= \frac{r_0 - n a_0}{p}
		= \frac{m - n}{3}
	.\]
	וכן,
	\[
		r_1 \equiv n a_1 \mod 3
		\iff r_1 \equiv 2 n \mod 3
	.\]
	ואף,
	\[
		m
		= r_0
		= r_2
		= \frac{r_1 - n a_1}{3}
		= \frac{r_1 - 2n}{3}
	.\]
	ולכן $3m = r_1 - 2n$ וכן,
	\[
		3m = \frac{m - n}{3} - 2n
		\iff 9m = m - n - 6n
		\iff 8m = -7n
	.\]
	אז,
	\[
		\frac{m}{n}
		= -\frac{7}{8}
	.\]
\end{solution}

\question{}
נמצא את הפיתוח ה־$3$־אדי הקנוני של $\frac{1}{4} \in \QQ_3$.
\begin{solution}
	מתקיים $|\frac{1}{4}|_3 = 1$ ולכן $s = 0$ וכן $m = 1, n = 4$ ונגדיר $\frac{1}{4} = \sum_{n = 0}^\infty a_n 3^n$.
	אז,
	\[
		r_0
		= m
		= 1
	.\]
	מהגדרת האיבר הראשון,
	\[
		1 \equiv 4 a_0 \mod 3
	.\]
	ולכן $a_0 = 1$ בדיוק.
	מתקיים,
	\[
		r_1
		= \frac{r_0 - n a_0}{p}
		= \frac{1 - 4 \cdot 1}{3}
		= -1
	.\]
	ולכן,
	\[
		r_1 \equiv n a_1 \mod 3
		\iff -1 \equiv 4 a_1 \mod 3
		\iff a_1 = 2
	.\]
	ואז מתקבל,
	\[
		r_2
		= \frac{r_1 - n a_1}{p}
		= \frac{-1 - 4 \cdot 2}{3}
		= -3
	.\]
	שוב,
	\[
		r_2 \equiv n a_2 \mod 3
		\iff -3 \equiv 4 a_2 \mod 3
		\iff a_2 = 0
	.\]
	ונקבל את השארית,
	\[
		r_3
		= \frac{-3 - 4 \cdot 0}{3}
		= -1
	.\]
	כלומר $r_1 = r_3$ ונסיק כי יש מחזור מגודל $2$ וכן,
	\[
		\frac{1}{4}
		= {(\ldots 0202021)}_3
	.\]
\end{solution}

\question{}
יהי $\FF((X^{-1}))$ השדה שהגדרנו בשאלה 1. \\
נמצא סדרה ${(a_n)}_{n = 0}^{\infty} \subseteq \FF$ כך שמתקיים,
\[
	\frac{1}{X - 1} = \sum_{n = 0}^\infty \frac{a^n}{X^n}
.\]
\begin{solution}
	נבחין כי $|X^{-n}|_{\alpha} = \alpha^{-(-n)} \xrightarrow{n \to \infty} 0$ ולכן $X^{-n} - 1 \to -1$.
	מתקיים גם,
	\[
		X^n - 1
		= (X - 1)(1 + \cdots + X^{n - 1})
		\iff 1 - X^{-n}
		= (X - 1)(1 + \cdots + X^{-n + 1})
	.\]
	כלומר,
	\[
		\frac{1}{X - 1}
		= \sum_{n = 0}^\infty X^{-n}
	.\]
	וקיבלנו $a_n = 1$ לכל $n \in \NN$.
\end{solution}

\question{}
נבדוק אם לפולינום $f = x^2 + x + 1$ יש שורשים ב־$K$.

\subquestion{}
כאשר $K = \QQ_p$ ו־$p = 5$.
\begin{solution}
	נשים לב כי לכל $z \in \ZZ_{/ 5}$ מתקיים,
	\[
		\overline{f(z)} \ne 0
	.\]
	ולכן אין שורש לפולינום.
\end{solution}

\subquestion{}
כאשר $p = 7$.
\begin{solution}
	אם קיים $w_0 \in \ZZ_p$ כך שמתקיים,
	\[
		\overline{f(w_0)} = 0,
		\qquad
		\overline{f'(w_0)} \ne 0
	.\]
	אז מהלמה של הנזל יש לפולינום שורש.
	נחשב את הנגזרת הפורמלית,
	\[
		f'
		= 2x + 1
	.\]

	נגדיר $w_0 = 2$ ונקבל,
	\[
		f(w_0)
		= 4 + 2 + 1
		\implies \overline{f(w_0)} = 0
	.\]
	וכן,
	\[
		f'(w_0)
		= 2 \cdot 2 + 1
		\implies \overline{f'(w_0)} \ne 0
	.\]
	ולכן קיים שורש לפולינום ב־$\QQ_7$ מהלמה של הנזל.
\end{solution}

\question{}
נמצא פולינום $f \in \ZZ_3[X]$ ללא שורשים ב־$\ZZ_3$ כך שלרדוקציה שלו ב־$\FF_3[X]$ יש שורשים ב־$\FF_3$.
\begin{solution}
	נגדיר את,
	\[
		f(X)
		= X^2 - 3
	.\]
	הצמצום שלו $f_0 \in \FF_3[X]$ הוא $f_0(X) = X^2$ ו־$X = 0$ שורש שלו, נשאר להראות שאין ל־$f$ עצמו שורש. \\
	נניח ש־$X \in \ZZ_3$ שורש של $f$. \\
	מהצד השני נניח בשלילה ש־$x \in \ZZ_3$ שורש של $f$, כלומר $x^2 = 3$.
	נקבל,
	\[
		|x|_3 \cdot |x|_3
		= |3|_p
		= \frac{1}{3}
	.\]
	כלומר $|x|_3 \notin \{ \frac{1}{3^n} \mid n \in \NN \}$ בסתירה למסקנות ממטלות קודמות.
	נובע של־$f$ אין שורשים ב־$\ZZ_3$.
\end{solution}

\question{}
נראה של־$f = x^2 - 2$ יש שני שורשים ב־$\QQ_7$ ונמצא את שלוש הספרות האחרונות בפיתוח ה־$7$־אדי הקנוני שלהם.
\begin{proof}
	נסמן $w_0 = 3$ ונקבל $f(w_0) \equiv 0 \mod 7$, וכן $f' = 2x$ ולכן גם $f'(w_0) \not\equiv 0 \mod 7$ ומהלמה של הנזל נובע שיש ל־$f$ שורש ב־$\QQ_7$.
	נסמן את השורש ב־$a$, ונשים לב ש־$-a \in \QQ_7$ אבל ${(-a)}^2 = a^2$ ולכן מצאנו שורש שני, ונסיק שהפולינום מתפרק לגורמים לינאריים מעל השדה, בפרט אין לו שורשים נוספים.

	נעבור לחישוב הספרות, נעשה זאת בהתבסס על האלגוריתם המוצג בהוכחת הלמה של הנזל.
	נניח ש־${(\ldots a_3 a_2 a_1)}_7 = a$ ונסמן את הסכום החלקי על־ידי $r_k$.
	נגדיר את $a_1 = 3$ האיבר היחיד $0 \le a_1 \le 6$ המקיים $\overline{a_1} = \overline{w_0}$.
	נסמן $z \in \ZZ_p$ עבורו $f(a_1) = z 7^1$, כלומר,
	\[
		f(3) = z \cdot 7
		\iff z = 1
	.\]
	ונגדיר את $a_2$ על־ידי הערך $0 \le a_2 \le 6$ היחיד המקיים,
	\[
		6 a_2
		= f'(a_1) a_2
		\equiv_7 -z
		= -1
	.\]
	נקבל ש־$a_2 = 1$.
	נגדיר שוב את,
	\[
		f(r_2)
		= z 7^2
		\implies f(3 + 1 \cdot 7)
		= z 7^2
		\implies 98 = 7^2 z
		\implies z = 2
	\]
	ונקבל,
	\[
		20 a_3
		= f'(r_2) a_3
		\equiv_7 -z
		= -2
		\implies 6a_3 \equiv_7 5
	.\]
	ונסיק ש־$a_3 = 2$.
	מצאנו אם כך שמתקיים,
	\[
		a = {(\ldots 213)}_7
	.\]

	נוכל עתה למצוא בקלות את הייצוג של $-a$, על־ידי חישוב ישיר,
	\[
		-a + a = 0
		\implies (-a) + {(\ldots 213)}_7
		= {(\ldots 000)}_7
		\implies -a
		= {(\ldots 454)}_7
	.\]
\end{proof}

\question{}
יהי $p$ ראשוני ויהיו $a, b, c \in p \ZZ_p$ כך ש־$c \notin p^2 \ZZ_p$. \\
נראה שלפולינום $f(x) = x^3 + a x^2 + b x + c$ אין שורשים ב־$\QQ_p$.
\begin{proof}
	TODO
\end{proof}

\end{document}
